% !TeX root = ../../Book1.tex
\chapter{弱拓扑、弱-*\WeakStarJoin{}拓扑与紧性}
\label{b1:ch:09}
\BookChapterNumber{b1:ch:09}
% 资料组：Fremlin2016Measure，2A5B--J、3A3I--J、3A5E--G、4A2Bd、4A4E、462A--D。
% 已读有关正文；3A3J、3A5F--G为结果与文献定位，不误称其附有本书所用证明。
% Xu2017Fanhan第8--9章仅核定官方目录、中文异名与伴读范围。
% 证明义务：有限测试与一般拓扑；乘积紧性；对偶球闭条件；Goldstine有限维分离；
% 自反性到可分子空间中的序列结论；凸分离到弱下半连续；极小化列的有界性。
\begin{chapterabstract}
在无穷维空间里，范数有界不能保证存在范数收敛子列。
然而，许多方程和极小化问题并不要求所有误差都以范数趋于零：
只要连续线性测试的结果收敛，就可能先保住方程，再借额外估计恢复更强的性质。
本章从这种测试方式构造弱拓扑和弱-*\WeakStarJoin{}拓扑，证明对偶球的紧性，
说明自反性如何把紧性带回原空间，最后把这些工具用于凸集、下半连续泛函和极小元的存在。
贯穿全章的问题是：减弱收敛以后，哪些信息仍然保存，哪些信息必须另行补回？
\end{chapterabstract}

% 本章目标（不计入编号）
\begin{learninggoals}
\begin{itemize}
\item\label{b1:item:09-goal-testing}区分范数收敛、弱收敛与弱-*\WeakStarJoin{}收敛，能根据所在空间写出正确的测试对象。
\item\label{b1:item:09-goal-compact}理解 Banach--Alaoglu 证明中的乘积空间与闭条件，知道紧性何时能够转化为收敛子列。
\item\label{b1:item:09-goal-reflexive}从自然嵌入识别自反空间，并能对自反空间中的有界列抽取弱收敛子列。
\item\label{b1:item:09-goal-convex}用凸分离和 Mazur 引理比较强、弱闭包，说明凸性在这种比较中的作用。
\item\label{b1:item:09-goal-minimum}将极小化列、有界性、弱闭约束及弱下半连续性组合成完整的存在性证明。
\end{itemize}
\end{learninggoals}

本章依赖前一章的连续对偶、Hahn--Banach 分离、一致有界原理和 Hilbert 空间表示定理。
初读时不必预先掌握一般局部凸空间理论；弱邻域的有限测试形式会直接给出。
一般拓扑中的紧性与序列紧性需要认真区别，为此正文补入网的闭包判据及乘积紧性证明。
可分空间中的序列版本则有可直接操作的对角抽取过程。

\ProseParagraphBreak

建议先在 \(\ell^2\)、\(c_0\) 和 \(\ell^1\) 上复算两节中的收敛例子，
然后连读\secref{b1:sec:0903}与\secref{b1:sec:0904}。
最后一节的变分应用只需要前面已证的自反空间序列结论，
不依赖选读部分的一般 Eberlein--Šmulian 定理。
这一安排使读者可以先掌握常用的存在性方法，再回头理解一般弱紧性的拓扑特征。

\ProseParagraphBreak

本章采用的定义和一般结果可配合 Fremlin 的《Measure Theory》%
\cite[Appendix to Volume 2, \S2A5; Appendix to Volume 3, \S3A5]{Fremlin2016Measure}阅读；
其附录按供后文调用的方式紧凑排列，本书则把主要证明与例子接成一条教学主线。
中文伴读可选许全华、马涛、尹智《泛函分析讲义》%
\cite[第8--9章]{Xu2017Fanhan}，对照其中的弱拓扑、弱-*\WeakStarJoin{}紧性及自反性各节。
希望进一步系统学习局部凸空间或变分方法，可继续阅读 Rudin 的《Functional Analysis》%
\cite{Rudin1991Functional}及 Brézis 的《Functional Analysis, Sobolev Spaces and Partial Differential Equations》%
\cite{Brezis2011Functional}；后者也可作为下一章至 Sobolev 空间部分的连续伴读。

% !TeX root = ../../Book1.tex
\section{弱拓扑}
\label{b1:sec:0901}

% 来源：Fremlin2016Measure，第2卷附录2A，2A5B、2A5D、2A5I，附录第24--27页。
% 实际阅读作者公开分章PDF；将半范数生成拓扑改写成有限测试邻域。
% 弱列的有界性、c0例子及网的闭包判别在此补出证明，不归称为原书的原有陈述。
范数收敛要求误差向量本身变小。由上一章的对偶范数公式，这也等于要求单位球内所有连续线性泛函读出的误差一致变小。然而，在许多极限问题中，我们只能对每个固定泛函分别证明收敛。弱拓扑就是为这种较少的控制量建立的拓扑。本节中 \(X\) 总是标量域 \(\mathbb K=\mathbb R\) 或 \(\mathbb C\) 上的赋范空间，暂不要求完备。

\subsection{弱拓扑}
\label{b1:subsec:090101}

\begin{definition}\label{b1:def:weak-topology}
对 \(x\in X\)、有限集 \(F\subset X^*\) 及 \(\varepsilon>0\)，令
\[
 U(x;F,\varepsilon)
 =\{\,y\in X\mid |f(y-x)|<\varepsilon\text{ 对每个 }f\in F\,\}.
\]
以这些集合为邻域基得到的拓扑称为 \(X\) 的\term[ruo-tuopu]{弱拓扑}[weak topology]，记为 \(\sigma(X,X^*)\)。具体说，集合 \(O\subset X\) 弱开，是指每个 \(x\in O\) 都有某个 \(U(x;F,\varepsilon)\subset O\)。
\end{definition}
\symidx{\(\sigma(X,X^*)\)：赋范空间的弱拓扑}

一个弱邻域只列出有限项测试，但测试泛函可以因邻域而异。取 \(F=\varnothing\) 时约定 \(U(x;F,\varepsilon)=X\)。不同测试使用不同的正误差容许量也给出同一拓扑：把有限个容许量取最小值，即可找到上述形式的更小邻域。

\begin{proposition}\label{b1:prop:weak-topology-properties}
弱拓扑是使所有 \(f\in X^*\) 连续的最弱拓扑；它是 Hausdorff 拓扑，且不强于范数拓扑。
\end{proposition}
\begin{proof}
先验证邻域定义。有限个同心基本邻域的交包含把测试集合合并、误差取最小值得到的基本邻域。若 \(y\in U(x;F,\varepsilon)\) 且 \(F\ne\varnothing\)，取
\[
 0<\delta<\min_{f\in F}\bigl(\varepsilon-|f(y-x)|\bigr),
\]
三角不等式给出 \(U(y;F,\delta)\subset U(x;F,\varepsilon)\)。空测试集合的情形显然。因此这些集合确实构成所定义拓扑的开邻域基。

固定 \(f\in X^*\)，条件 \(|f(y)-f(x)|<\varepsilon\) 本身就是一个基本邻域条件，故 \(f\) 弱连续。反过来，任何使全部 \(f\) 连续的拓扑，都使这些标量开球反像的有限交开，从而包含弱拓扑。这证明最弱性。

若 \(x\ne y\)，由\cref{b1:cor:norming-functional}可取 \(f\in X^*\) 使 \(d=|f(x-y)|>0\)。邻域 \(U(x;\{f\},d/3)\) 与 \(U(y;\{f\},d/3)\) 不相交，否则三角不等式会给出 \(d<2d/3\)。最后，有限个估计
\[
 |f(y-x)|\leq\|f\|\,\|y-x\|
\]
保证每个基本弱邻域都包含同心的范数开球，故每个弱开集也是范数开集。
\end{proof}

加法与数乘在弱拓扑下连续：用任意 \(f\in X^*\) 测试后，它们分别变成标量的加法和数乘；有限项测试的误差可以同时控制。弱连续的线性泛函仍恰好是 \(X^*\)，因为弱连续蕴含范数连续，而反向包含已经由定义保证。

有限维时，选一组基，其坐标泛函均连续；有限个坐标同时变小，就由\cref{b1:prop:finite-dimensional-norms}中的估计迫使范数变小，所以弱拓扑与范数拓扑相同。无穷维时则严格不同。给定有限集 \(F=\{f_1,\ldots,f_m\}\)，线性映射 \(z\mapsto(f_1(z),\ldots,f_m(z))\) 不可能从无穷维空间单射到 \(\mathbb K^m\)，故有非零 \(v\in\bigcap_{f\in F}\ker f\)。整条直线 \(x+\mathbb Rv\) 都包含在 \(U(x;F,\varepsilon)\) 内。因此每个弱邻域都在范数下无界，范数开单位球不可能弱开。

\subsection{弱收敛}
\label{b1:subsec:090102}

\begin{definition}\label{b1:def:weak-convergence}
序列 \((x_n)\subset X\) 称为\term[ruo-shoulian]{弱收敛}[weak convergence]于 \(x\in X\)，记为 \(x_n\rightharpoonup x\)，若
\[
 f(x_n)\longrightarrow f(x)\quad\text{对每个 }f\in X^*.
\]
\end{definition}
\symidx{\(x_n\rightharpoonup x\)：序列的弱收敛}

这正是序列在弱拓扑下的收敛。逐个泛函收敛时，对一个基本邻域中的有限个泛函分别选取起始指标，再取最大值，序列便最终落入该邻域；反向结论只需使用单个泛函的邻域。弱极限由 Hausdorff 性唯一确定。若 \(T\in\mathcal B(X,Y)\)，则 \(x_n\rightharpoonup x\) 蕴含 \(Tx_n\rightharpoonup Tx\)，因为对任意 \(g\in Y^*\)，复合 \(gT\) 属于 \(X^*\)。

逐项测试的收敛并不直接给出一个对全部测试统一的误差估计，却仍会给出原序列的范数界。

\begin{proposition}\label{b1:prop:weak-sequence-bounded}
若序列 \((x_n)\subset X\) 满足：对每个 \(f\in X^*\)，标量列 \((f(x_n))\) 有界，则 \(\sup_n\|x_n\|<\infty\)。特别地，每个弱收敛列都在范数下有界。这里不要求 \(X\) 完备。
\end{proposition}
\begin{proof}
把 \(x_n\) 看成 \(X^*\) 上的泛函 \(J_Xx_n\)，其中 \((J_Xx_n)(f)=f(x_n)\)。由\cref{b1:prop:canonical-embedding}，\(\|J_Xx_n\|=\|x_n\|\)。又由\cref{b1:thm:operator-space-complete}，\(X^*\) 总是 Banach 空间。题设正是算子族 \(\{J_Xx_n\}\) 在 \(X^*\) 上逐点有界，故\cref{b1:thm:uniform-boundedness}给出 \(\sup_n\|J_Xx_n\|<\infty\)。
\end{proof}

本书讨论弱收敛时，首先讨论上述序列。不过，一般拓扑中的闭包和紧致性不能未经证明就用序列代替。为保留必要的区别，我们现在引入一种允许更一般指标的收敛。

\begin{definition}\label{b1:def:net}
非空偏序集 \(D\) 称为\term[youxiang-ji]{有向集}[directed set]，若任意两个元素都有共同的上界。以 \(D\) 为指标集的点族 \((x_d)_{d\in D}\) 称为\term[wang]{网}[net]。在拓扑空间中，称 \(x_d\to x\)，若对 \(x\) 的每个邻域 \(V\)，都存在 \(d_0\in D\)，使 \(d\geq d_0\) 时 \(x_d\in V\)。
\end{definition}

自然数的通常次序是一个有向集，所以序列是网的特例。有向性保证有限多个起始指标有共同上界；因此，网在弱拓扑中收敛仍等价于对每个 \(f\in X^*\) 都有 \(f(x_d)\to f(x)\)。

\begin{proposition}\label{b1:prop:net-closure}
设 \(A\) 是拓扑空间 \(Z\) 的子集。点 \(x\) 属于 \(A\) 的闭包，当且仅当存在一个取值于 \(A\) 的网收敛到 \(x\)。
\end{proposition}
\begin{proof}
若有这样的网，\(x\) 的每个邻域都包含网中的点，因而与 \(A\) 相交。反过来，若 \(x\in\overline A\)，令 \(D\) 为 \(x\) 的全部开邻域，按反包含关系排序，即 \(U\leq V\) 表示 \(V\subset U\)。两邻域的交是共同上界，所以 \(D\) 有向。对每个 \(U\in D\) 选取 \(a_U\in A\cap U\)。给定 \(x\) 的开邻域 \(V\)，当 \(U\geq V\) 时有 \(a_U\in U\subset V\)，故 \(a_U\to x\)。
\end{proof}

度量空间中可以只取半径 \(1/n\) 的球，由此得到闭包的序列判别。弱拓扑未必有这样的可数邻域基，所以上述命题一般必须保留网。还应注意：收敛网的一个起始部分未必有限，不能照搬“收敛序列有界”的论证，把\cref{b1:prop:weak-sequence-bounded}擅自改成关于所有弱收敛网的断言。

\subsection{弱 Cauchy 序列}
\label{b1:subsec:090103}

\begin{definition}\label{b1:def:weak-cauchy}
若对每个 \(f\in X^*\)，标量列 \((f(x_n))\) 都是 Cauchy 列，则称 \((x_n)\) 为\term[ruo-cauchy-xulie]{弱 Cauchy 序列}[weakly Cauchy sequence]。
\end{definition}

弱收敛列必为弱 Cauchy 序列。反过来，标量域完备保证每个测试都产生一个极限，但这些极限能否同时由 \(X\) 中的同一个向量产生，是另一个问题。先由\cref{b1:prop:weak-sequence-bounded}取 \(M=\sup_n\|x_n\|<\infty\)，再定义
\[
 \Lambda(f)=\lim_{n\to\infty}f(x_n)\quad(f\in X^*).
\]
标量极限与线性组合交换，且 \(|\Lambda(f)|\leq M\|f\|\)，故 \(\Lambda\in X^{**}\)。原序列在 \(X\) 中弱收敛，当且仅当这个泛函可写成 \(J_Xx\)。这说明障碍并非标量极限不存在，而是极限可能无法在原空间中表示。

即使 \(X\) 是 Banach 空间，也可能出现这一障碍。取配备一致范数的 \(c_0\)，并令
\[
 s_n=(\underbrace{1,\ldots,1}_{n\text{ 项}},0,0,\ldots).
\]
由\cref{b1:ex:08-c0-dual}，每个 \(f\in c_0^*\) 唯一写成 \(f(x)=\sum_{k\geq1}a_kx_k\)，其中 \(a\in\ell^1\)。于是
\[
 f(s_n)=\sum_{k=1}^n a_k\longrightarrow\sum_{k=1}^{\infty}a_k,
\]
所以 \((s_n)\) 弱 Cauchy。若它弱收敛于 \(x\in c_0\)，逐个使用坐标泛函就得到 \(x_k=1\) 对所有 \(k\) 成立，与 \(x\in c_0\) 矛盾。范数完备性只保证范数 Cauchy 列的极限，不保证这里的较弱条件。

\subsection{范数收敛与弱收敛}
\label{b1:subsec:090104}

由 \(|f(x_n-x)|\leq\|f\|\,\|x_n-x\|\)，范数收敛必然蕴含弱收敛；为突出区别，范数收敛也称为\term[qiang-shoulian]{强收敛}[strong convergence]。无穷维空间中，反向蕴含一般失败，而且失败可以发生在最熟悉的 Hilbert 空间里。

设 \((e_n)\) 是 Hilbert 空间 \(H\) 中的正交规范列。对任意 \(h\in H\)，由\cref{b1:prop:bessel-inequality}有
\[
 \sum_{n=1}^{\infty}|\langle e_n,h\rangle|^2\leq\|h\|^2,
\]
因此 \(\langle e_n,h\rangle\to0\)。再由\cref{b1:thm:hilbert-riesz}，全部连续线性泛函都能如此表示，所以 \(e_n\rightharpoonup0\)，但 \(\|e_n\|=1\)。弱收敛允许向量的方向不断改变，使任何一个固定测试最终都看不到它，而并未要求向量长度消失。

\begin{proposition}\label{b1:prop:hilbert-weak-norm-convergence}
在 Hilbert 空间中，\(x_n\to x\) 强收敛，当且仅当 \(x_n\rightharpoonup x\) 且 \(\|x_n\|\to\|x\|\)。
\end{proposition}
\begin{proof}
正向结论由前述估计及范数的连续性成立。反过来，弱收敛给出 \(\langle x_n,x\rangle\to\langle x,x\rangle\)，故
\[
 \|x_n-x\|^2
 =\|x_n\|^2+\|x\|^2-2\operatorname{Re}\langle x_n,x\rangle
 \longrightarrow0.
\]
\end{proof}

这个结论利用了内积恒等式，不能无条件推广到一般赋范空间。在 \(c_0\) 中取 \(y_n=e_1+e_n\)（\(n\geq2\)）。仍用 \(c_0^*=\ell^1\) 的表示，因为 \(a_n\to0\)，有 \(f_a(y_n)=a_1+a_n\to a_1=f_a(e_1)\)，故 \(y_n\rightharpoonup e_1\)。尽管 \(\|y_n\|_\infty=\|e_1\|_\infty=1\)，距离 \(\|y_n-e_1\|_\infty\) 始终等于 \(1\)。因此，弱收敛、范数有界和范数值的收敛是应当分别核对的三件事。

% !TeX root = ../../Book1.tex
\section{弱-*\WeakStarJoin{}拓扑}
\label{b1:sec:0902}

% 来源：Fremlin2016Measure，第2卷附录2A5I(g)，附录第27页，实际阅读作者公开PDF。
% 弱-*与弱拓扑的有限测试解释据此改写；l1对偶、坐标判别与尾部示性序列在此完整推导。
当向量本身就是泛函时，有一种更直接的测试办法：给定一个输入，观察泛函在该输入上的值。这样的测试只使用原空间的向量，而不使用对偶空间上的全部连续线性泛函，因而产生了另一种拓扑。弱拓扑与弱-*\WeakStarJoin{}拓扑的区别，首先是允许使用的测试族不同。

\subsection{对偶空间上的弱-*\WeakStarJoin{}拓扑}
\label{b1:subsec:090201}

\begin{definition}\label{b1:def:weak-star-topology}
设 \(X\) 为赋范空间。对 \(f\in X^*\)、有限集 \(E\subset X\) 与 \(\varepsilon>0\)，令
\[
 V(f;E,\varepsilon)
 \coloneqq\{\,g\in X^*\mid |g(x)-f(x)|<\varepsilon\text{ 对每个 }x\in E\,\}.
\]
以这些集合为邻域基的拓扑称为 \(X^*\) 上的\term[ruo-xing-tuopu]{弱-*\WeakStarJoin{}拓扑}[weak-* topology]，记为 \(\sigma(X^*,X)\)。
\end{definition}
\symidx{\(\sigma(X^*,X)\)：对偶空间相对于原空间的弱-*\WeakStarJoin{}拓扑}

与\cref{b1:prop:weak-topology-properties}中的有限测试论证完全相同，这些集合组成开邻域基，而 \(\sigma(X^*,X)\) 是使所有求值映射
\[
 X^*\longrightarrow\mathbb K,\quad f\longmapsto f(x)\quad(x\in X)
\]
连续的最弱拓扑。若 \(f\ne g\)，总有一个 \(x\in X\) 使 \(f(x)\ne g(x)\)，故两个小求值邻域可将它们分开，弱-*\WeakStarJoin{}拓扑也是 Hausdorff 的。估计 \(|(g-f)(x)|\leq\|g-f\|\cdot\|x\|\) 又说明范数拓扑不弱于弱-*\WeakStarJoin{}拓扑。这些事实都不要求 \(X\) 完备。

\ProseParagraphBreak

记号中保留 \(X\) 十分重要。这里不是任意给 \(X^*\) 降低一个拓扑，而是明确规定只用来自 \(X\) 的求值泛函作为测试。Fremlin 的《Measure Theory》\cite[Appendix to Volume 2, 2A5I]{Fremlin2016Measure}把弱拓扑与弱-*\WeakStarJoin{}拓扑放在同一框架下讨论，适合用来对照这两种测试族；下面则用序列空间把区别具体化。

\subsection{弱-*\WeakStarJoin{}收敛}
\label{b1:subsec:090202}

\begin{definition}\label{b1:def:weak-star-convergence}
设 \((f_n)\subset X^*\) 且 \(f\in X^*\)。若 \(f_n(x)\to f(x)\) 对每个 \(x\in X\) 都成立，则称它\term[ruo-xing-shoulian]{弱-*\WeakStarJoin{}收敛}[weak-* convergence]于 \(f\)，记为 \(f_n\overset{*}{\rightharpoonup}f\)。
\end{definition}
\symidx{\(f_n\overset{*}{\rightharpoonup}f\)：泛函序列的弱-*\WeakStarJoin{}收敛}

这就是对偶空间内部的逐点收敛。条件 \(f\in X^*\) 是定义的一部分，不能只从逐点极限公式中省略。事实上，在赋予一致范数的 \(c_{00}\) 上令
\[
 f_n(x)=\sum_{k=1}^n x_k.
\]
每个 \(f_n\) 都连续；对每个有限支撑的 \(x\)，它逐点趋于 \(\sum_{k\geq1}x_k\)，但这个极限泛函在一致范数下不连续，因为它在 \(x^{(N)}=\sum_{k=1}^Ne_k\) 上取值 \(N\)，而 \(\|x^{(N)}\|_\infty=1\)。所以当原空间不完备时，连续泛函的逐点极限未必仍是弱-*\WeakStarJoin{}拓扑中的点。

\ProseParagraphBreak

有限多个输入对应的起始指标取最大值，便得到序列在弱-*\WeakStarJoin{}拓扑中的收敛；网也有同样的逐点判别，只须把最大值换成有限个指标的共同上界。以下有界性命题则仍明确针对序列。

\begin{proposition}\label{b1:prop:weak-star-sequence-bounded}
若 \(X\) 是 Banach 空间，且 \(f_n\overset{*}{\rightharpoonup}f\) 于 \(X^*\)，则 \(\sup_n\|f_n\|<\infty\)。
\end{proposition}
\begin{proof}
对每个 \(x\in X\)，收敛的标量列 \(\bigl(f_n(x)\bigr)\) 有界。把 \(f_n\) 看成从 Banach 空间 \(X\) 到 \(\mathbb K\) 的有界线性算子，\cref{b1:thm:uniform-boundedness}直接给出结论。
\end{proof}

这里的一致有界原理作用在 \(X\) 上，而\cref{b1:prop:weak-sequence-bounded}中作用在自动完备的 \(X^*\) 上，故两条命题的假设不同。\subsecref{b1:subsec:080302}中的例子恰好说明完备性不能直接删去：在赋予一致范数的 \(c_{00}\) 上，令 \(f_n(x)=nx_n\)。每个固定输入只有有限个非零坐标，所以 \(f_n(x)\) 最终等于零，即 \(f_n\overset{*}{\rightharpoonup}0\)；但 \(\|f_n\|=n\)。对偶空间 \(c_{00}^*\) 本身完备，并不能补上定义域的这个缺口。

\ProseParagraphBreak

为了把弱-*\WeakStarJoin{}收敛写成常见的坐标条件，我们先确定一个具体对偶空间。

\begin{proposition}[\(\ell^1\) 的完备性与连续对偶]\label{b1:prop:l1-dual}
空间 \(\ell^1\) 关于范数 \(\|a\|_1=\sum_k|a_k|\) 完备。对 \(u\in\ell^\infty\)，公式
\[
 F_u(a)=\sum_{k=1}^{\infty}a_ku_k\quad(a\in\ell^1)
\]
定义 \((\ell^1)^*\) 中的元素。映射 \(u\mapsto F_u\) 是 \(\ell^\infty\) 到 \((\ell^1)^*\) 的线性等距满射。
\end{proposition}
\begin{proof}
先设 \((a^{(j)})\) 是 \(\ell^1\) 中的 Cauchy 列。每个坐标列都是 Cauchy 列，记 \(a_k=\lim_j a_k^{(j)}\)。给定 \(\varepsilon>0\)，取 \(J\) 使 \(j,m\geq J\) 时 \(\|a^{(j)}-a^{(m)}\|_1<\varepsilon\)。固定 \(j\geq J\)，对任意 \(N\) 令 \(m\to\infty\)，得到
\[
 \sum_{k=1}^N|a_k^{(j)}-a_k|\leq\varepsilon.
\]
令 \(N\to\infty\)，可知 \(a-a^{(j)}\in\ell^1\) 且 \(\|a-a^{(j)}\|_1\leq\varepsilon\)。因此 \(a\in\ell^1\) 且 \(a^{(j)}\to a\)，完备性得证。

绝对收敛性和估计 \(|F_u(a)|\leq\|u\|_\infty\cdot\|a\|_1\) 给出 \(\|F_u\|\leq\|u\|_\infty\)。逐个代入单位坐标向量 \(e_k\)，得到 \(\|F_u\|\geq|u_k|\)，取上确界即有范数相等。

反过来，给定 \(F\in(\ell^1)^*\)，令 \(u_k=F(e_k)\)，则 \(|u_k|\leq\|F\|\)，所以 \(u\in\ell^\infty\)。对 \(a\in\ell^1\)，其有限坐标截断 \(a^{(N)}\) 在 \(\ell^1\) 范数下趋于 \(a\)，故
\[
 F(a)=\lim_{N\to\infty}F(a^{(N)})
 =\lim_{N\to\infty}\sum_{k=1}^N a_ku_k=F_u(a).
\]
各坐标 \(u_k\) 已由 \(F(e_k)\) 唯一确定，因而表示唯一。
\end{proof}

上述配对在实、复情形下都不插入共轭；这样，\(a\mapsto F_u(a)\) 与 \(u\mapsto F_u\) 都是线性的。它不同于 Hilbert 空间的内积表示约定。

\begin{corollary}\label{b1:cor:linfty-weak-star-coordinate}
在识别 \(\ell^\infty=(\ell^1)^*\) 后，序列 \((u^{(n)})\) 弱-*\WeakStarJoin{}收敛于 \(u\)，当且仅当
\[
 \sup_n\|u^{(n)}\|_\infty<\infty,
 \quad u_k^{(n)}\longrightarrow u_k\quad\text{对每个 }k.
\]
\end{corollary}
\begin{proof}
弱-*\WeakStarJoin{}收敛时，用 \(e_k\in\ell^1\) 测试便得到坐标收敛。\Cref{b1:prop:l1-dual}已经直接证明 \(\ell^1\) 是 Banach 空间，故\cref{b1:prop:weak-star-sequence-bounded}给出统一界。

反过来，设上述统一界为 \(M\)。坐标极限满足 \(|u_k|\leq M\)，所以 \(u\in\ell^\infty\)。对任意 \(a\in\ell^1\)，分成前 \(N\) 项与尾部可得
\[
 \left|\sum_{k=1}^{\infty}a_k(u_k^{(n)}-u_k)\right|
 \leq\sum_{k=1}^N|a_k|\cdot|u_k^{(n)}-u_k|
      +2M\sum_{k>N}|a_k|.
\]
先选 \(N\) 使尾部任意小，再由有限个坐标的收敛使首项任意小，即得到每个 \(a\) 对应的求值收敛。
\end{proof}

有限个坐标测试本身看不到尾部；统一范数界与测试序列的绝对可和性一起，才把尾部误差控制住。这个先截断再取极限的步骤，也解释了为什么不能只凭坐标收敛就断言弱-*\WeakStarJoin{}收敛。

\subsection{弱拓扑与弱-*\WeakStarJoin{}拓扑的区别}
\label{b1:subsec:090203}

把 \(X^*\) 当作赋范空间，它自己的弱拓扑是 \(\sigma(X^*,X^{**})\)，允许使用所有 \(\Phi\in X^{**}\) 测试；弱-*\WeakStarJoin{}拓扑 \(\sigma(X^*,X)\) 只允许其中形如 \(J_Xx\) 的测试。因此在 \(X^*\) 上有
\[
 \text{范数拓扑}\ \supseteq\ \sigma(X^*,X^{**})
 \ \supseteq\ \sigma(X^*,X).
\]
这里的包含表示开集族的包含，相应地，范数收敛蕴含弱收敛，弱收敛蕴含弱-*\WeakStarJoin{}收敛。如果 \(J_X(X)=X^{**}\)，后两种拓扑便相同；一般情况下，缺少的测试确实可能改变收敛性。

\ProseParagraphBreak

自然嵌入本身还有一个拓扑性质，后面把原空间放进二次对偶时会反复使用。

\begin{proposition}\label{b1:prop:canonical-weak-embedding}
自然嵌入 \(J_X\) 把 \(\bigl(X,\sigma(X,X^*)\bigr)\) 同胚地映到 \(X^{**}\) 中的像 \(J_X(X)\)，其中像空间采用 \(X^{**}\) 上弱-*\WeakStarJoin{}拓扑 \(\sigma(X^{**},X^*)\) 的子空间拓扑。
\end{proposition}
\begin{proof}
这里要说明 \(J_X\) 是单射。而 \(X^{**}\) 中以 \(J_Xx\) 为中心、由有限个 \(f_1,\ldots,f_m\in X^*\) 给出的弱-*\WeakStarJoin{}基本邻域，在 \(J_X\) 下的反像由
\[
 |(J_Xy-J_Xx)(f_j)|=|f_j(y-x)|<\varepsilon
 \quad(1\leq j\leq m)
\]
刻画。这恰好是 \(X\) 的基本弱邻域，且所有基本弱邻域都如此得到。因此 \(J_X\) 及其在像上的逆映射都连续。
\end{proof}

来看一个弱-*\WeakStarJoin{}收敛却不弱收敛的序列。在实空间 \(\ell^\infty=(\ell^1)^*\) 中，令
\[
 u_k^{(n)}=\mathbf1_{\{k\geq n\}}.
\]
它的范数恒等于 \(1\)，每个固定坐标都最终变成零，所以由\cref{b1:cor:linfty-weak-star-coordinate}有 \(u^{(n)}\overset{*}{\rightharpoonup}0\)。另一方面，令 \(c\subset\ell^\infty\) 为实收敛序列组成的子空间。通常极限映射 \(\lambda:c\to\mathbb R\) 是范数为一的连续线性泛函，由\cref{b1:thm:hahn-banach}可将它保范延拓为某个 \(L\in(\ell^\infty)^*\)。对每个固定 \(n\)，序列 \((u_k^{(n)})_{k\geq1}\) 的通常极限为 \(1\)，因此 \(L(u^{(n)})=1\)。这个合法的弱测试不趋于零，故 \((u^{(n)})\) 不弱收敛于零。它也不可能弱收敛于别的向量，因为弱收敛蕴含弱-*\WeakStarJoin{}收敛，而弱-*\WeakStarJoin{}极限唯一。

\ProseParagraphBreak

这里使用尾部示性序列是有意的：单位坐标向量 \(e_n\) 在 \(\ell^\infty\) 中反而弱收敛于零。事实上，对任意 \(L\in(\ell^\infty)^*\)，在每个有限和中选择模为 \(1\) 的标量 \(\alpha_k\)，使 \(\alpha_k\cdot L(e_k)=|L(e_k)|\)，便有
\[
 \sum_{k=1}^N|L(e_k)|
 =L\left(\sum_{k=1}^N\alpha_ke_k\right)\leq\|L\|.
\]
故 \(L(e_n)\to0\)。相同的坐标外观并不足以判断弱收敛，必须看清所处的空间及它允许的全部泛函。

\ProseParagraphBreak

弱-*\WeakStarJoin{}拓扑保留逐点测试，却舍去了一部分弱测试。下一节将说明，这种舍弃换来一个重要的紧致性结论；使用那个结论时，既要记住原空间是谁，也不能把一般拓扑的紧致性预先当作子列收敛性。

% !TeX root = ../../Book1.tex
\section{Banach–Alaoglu 定理}
\label{b1:sec:0903}

在范数拓扑中，无穷维空间的闭单位球不紧；换用弱-*拓扑以后，对偶单位球却总是紧的。
这个结论不要求原空间完备，也不要求可分。证明的起点很具体：
一个范数不超过一的泛函，在每个固定向量上的取值都落在一个紧圆盘内。
困难在于这些向量通常不可数，不能直接用可数次抽取子列处理。

\subsection{对偶单位球}
\label{b1:subsec:090301}

记 \(B_{X^*}=\{f\in X^*:\|f\|\leq1\}\)。对每个 \(x\in X\)，令
\[
 D_x=\{z\in\mathbb K:|z|\leq\|x\|\},
 \quad P=\prod_{x\in X}D_x.
\]
这里 \(\mathbb K=\mathbb R\) 或 \(\mathbb C\)；实情形的 \(D_x\) 是闭区间，复情形是闭圆盘。
把泛函送到它的全部取值，就得到单射
\[
 \Phi:B_{X^*}\longrightarrow P,
 \quad \Phi(f)=(f(x))_{x\in X}.
\]
乘积中的一般元素只是一份取值表，不一定满足线性关系。证明将先得到整份取值表空间的紧性，
再把线性关系作为闭条件加进去。

在乘积空间 \(\prod_{i\in I}K_i\) 上，\term[chengji-tuopu]{乘积拓扑}[product topology]的基本开集只限制有限个坐标：
\[
 \bigcap_{i\in F}\pi_i^{-1}(U_i),
 \quad F\subset I\text{ 有限},\quad U_i\subset K_i\text{ 开},
\]
其中 \(\pi_i\) 是第 \(i\) 个坐标投影。因此 \(\Phi\) 把弱-*基本邻域变成像空间中的乘积基本邻域。
这正是弱-*拓扑与逐点测试之间的关系，而不是一种另加的紧性假设。

本节暂时需要一般拓扑中的紧性：每个开覆盖有有限子覆盖。
取补集可知，这等价于每个具有\term[youxian-jiao-xingzhi]{有限交性质}[finite intersection property]的闭集族都有公共点；
有限交性质是指任意有限个成员的交非空。
下面给出乘积紧性所需的证明，以免把本章关键一步推给尚未展开的拓扑附录。
乘积拓扑和定理的一般版本可参阅 Fremlin\cite[卷三，3A3I--J]{Fremlin2016Measure}。

\begin{theorem}[Tychonoff]\label{b1:thm:tychonoff}
任意一族紧拓扑空间的乘积，在乘积拓扑下仍然紧。
\end{theorem}
\begin{proof}
若有因子为空，乘积为空，结论成立。以下设各因子 \(K_i\) 非空，记 \(P=\prod_iK_i\)。
本证明使用选择公理及其等价形式 Zorn 引理。
给定 \(P\) 中具有有限交性质的闭集族 \(\mathcal F\)，
把它扩充为具有有限交性质的极大子集族 \(\mathcal U\)。
这样的扩充存在，因为按包含关系排序的一条链的并仍具有有限交性质：
有限多个成员已同时属于链中的某一个族。

极大性给出三个性质。首先，\(\mathcal U\) 对有限交封闭，
因为加入已有成员的有限交不会破坏有限交性质。其次，若 \(A\in\mathcal U\) 且 \(A\subset B\subset P\)，
则 \(B\in\mathcal U\)，理由相同。最后，对任意 \(B\subset P\)，
\(B\) 与 \(P\setminus B\) 中恰有一个属于 \(\mathcal U\)。
若 \(B\notin\mathcal U\)，加入 \(B\) 会破坏有限交性质，因而存在有限多个已有成员，
其交 \(A\) 与 \(B\) 不交；前两个性质便推出 \(P\setminus B\in\mathcal U\)。
二者不能同时属于 \(\mathcal U\)，因为空集不属于这个族。

对每个 \(i\)，闭集族
\[
 \{\overline{\pi_i(A)}:A\in\mathcal U\}
\]
具有有限交性质；任取有限多个 \(A\) 的一个公共点并投影，就能验证这一点。
由 \(K_i\) 紧，可选
\(x_i\in\bigcap_{A\in\mathcal U}\overline{\pi_i(A)}\)，并令 \(x=(x_i)_{i\in I}\)。
若 \(U_i\) 是 \(x_i\) 的开邻域，则 \(\pi_i^{-1}(U_i)\in\mathcal U\)。
否则其补集属于 \(\mathcal U\)，而该补集的投影包含在闭集 \(K_i\setminus U_i\) 中，
与 \(x_i\) 的选取矛盾。

故 \(x\) 的每个乘积基本邻域都属于 \(\mathcal U\)，与每个 \(F\in\mathcal F\) 相交。
这说明 \(x\in\overline F=F\)。于是原闭集族有公共点，\(P\) 紧。
\end{proof}

证明中有限个坐标可以同时控制，是乘积拓扑的决定性特征。
若把“有限”改成“任意多个”，得到的拓扑一般更强，上述邻域论证便不再成立。

\subsection{Banach–Alaoglu 定理}
\label{b1:subsec:090302}

\begin{theorem}[Banach--Alaoglu]\label{b1:thm:banach-alaoglu}
设 \(X\) 是赋范空间。对偶闭单位球 \(B_{X^*}\) 在弱-*拓扑 \(\sigma(X^*,X)\) 下是紧的。
\end{theorem}
\begin{proof}
每个 \(D_x\) 紧，由\cref{b1:thm:tychonoff}，\(P\) 紧。
在 \(P\) 中考虑同时满足下列关系的取值表 \(u=(u_x)_{x\in X}\)：
\[
 u_{ax+by}=au_x+bu_y
 \quad(x,y\in X,\ a,b\in\mathbb K).
\]
对固定的 \(x,y,a,b\)，等式两边的差是连续的有限坐标函数，所以该等式定义闭集。
所有等式的公共解集因而是 \(P\) 的闭子集。

每张这样的表定义一个线性泛函 \(f(x)=u_x\)，且 \(|f(x)|\leq\|x\|\)，
故 \(f\in B_{X^*}\)。反过来，\(B_{X^*}\) 中每个泛函的取值表满足全部等式。
公共解集因此恰为 \(\Phi(B_{X^*})\)，它是紧集。
由于 \(\Phi\) 及其到像上的逆都把有限测试邻域对应起来，
\(\Phi\) 是弱-*拓扑到像的同胚，所求紧性成立。
\end{proof}

\term*[banach-alaoglu-dingli]{Banach--Alaoglu 定理}[Banach--Alaoglu theorem]中的星号不可省去。
它说的是 \(X^*\) 上由 \(X\) 测试的紧性，而不是由全部 \(X^{**}\) 测试的弱紧性。
闭球与有界性的条件也各有作用：对任意 \(R\geq0\)，半径为 \(R\) 的对偶闭球弱-*紧，
从而每个范数有界、弱-*闭的子集弱-*紧；一般范数闭集却不必弱-*闭。
Fremlin\cite[卷三，3A5F]{Fremlin2016Measure}采用同一赋范空间版本，亦不要求 \(X\) 完备。

下一节要把对偶球的紧性传回 \(X\)。为此，还需知道自然嵌入的像在二次对偶球中有多大。
它未必等于整个球，但有限多个泛函不能把球中的点与这个像分开。

\begin{theorem}[Goldstine]\label{b1:thm:goldstine}
设 \(J_X:X\to X^{**}\) 是自然等距嵌入。则
\[
 \overline{J_X(B_X)}^{\,\sigma(X^{**},X^*)}=B_{X^{**}}.
\]
\end{theorem}
\begin{proof}
右边的球由闭条件 \(|x^{**}(f)|\leq\|f\|\) 描述，故弱-*闭，并包含左边的闭包。
反过来，固定 \(x^{**}\in B_{X^{**}}\)、\(f_1,\ldots,f_m\in X^*\) 和 \(\varepsilon>0\)。
令
\[
 T:X\to\mathbb K^m,\quad Tx=(f_1(x),\ldots,f_m(x)),
 \quad z=(x^{**}(f_1),\ldots,x^{**}(f_m)).
\]
我们断言 \(z\in\overline{T(B_X)}\)。若不然，有限维空间中的凸分离定理给出一个实线性泛函，
把 \(z\) 与这个闭凸集严格分离。实情形它写成 \(w\mapsto\sum_i a_iw_i\)；
复情形则写成 \(w\mapsto\operatorname{Re}\sum_i a_iw_i\)。
记 \(g=\sum_i a_if_i\)，两种情形统一给出
\[
 \operatorname{Re}x^{**}(g)>
 \sup_{x\in B_X}\operatorname{Re}g(x)=\|g\|.
\]
最后的等号来自单位球的对称性；复情形还可乘以模为一的标量。
但 \(|x^{**}(g)|\leq\|g\|\)，矛盾。
于是能选 \(x\in B_X\)，使所有 \(|f_i(x)-x^{**}(f_i)|<\varepsilon\)。
这正是 \(x^{**}\) 属于所求弱-*闭包的邻域判据。
\end{proof}

\term*[goldstine-dingli]{Goldstine 定理}[Goldstine theorem]给出的是有限测试意义的逼近，
并不声称 \(J_X(B_X)\) 在二次对偶范数下稠密。
事实上，若 \(X\) 完备，等距像 \(J_X(X)\) 已是范数闭子空间；
除非自然嵌入满射，否则它不可能在 \(X^{**}\) 中范数稠密。

\subsection{可分情形的序列版本}
\label{b1:subsec:090303}

紧性在一般拓扑下不自动提供收敛子列。下面的可分性假设使有界对偶球只需可数个测试，
由此恢复度量空间中熟悉的子列结论。

\begin{theorem}\label{b1:thm:weak-star-metrizable}
设 \(X\) 是可分赋范空间，\(R>0\)。在 \(RB_{X^*}\) 上，弱-*拓扑可由一个度量给出。
特别地，每个范数有界的 \(X^*\) 中的序列都有弱-*收敛子列，其极限仍在 \(X^*\) 中。
\end{theorem}
\begin{proof}
取 \(X\) 的可数稠密集 \(\{x_j:j\geq1\}\)，在 \(RB_{X^*}\) 上定义
\[
 d(f,g)=\sum_{j=1}^{\infty}2^{-j}\min\{1,|(f-g)(x_j)|\}.
\]
对称性及三角不等式直接来自绝对值和
\(\min\{1,a+b\}\leq\min\{1,a\}+\min\{1,b\}\)。
若 \(d(f,g)=0\)，则 \(f,g\) 在稠密集上一致，由连续性知 \(f=g\)，所以 \(d\) 是度量。

先说明它确实给出弱-*拓扑。给定一个 \(d\)-球，先使级数尾和小于半个半径，
再限制前面有限多个测试值，就得到包含在球内的弱-*邻域。
反过来，考虑有限个条件 \(|(g-f)(y_i)|<\varepsilon\)。
分别选 \(x_{j_i}\) 逼近 \(y_i\)，使 \(2R\|y_i-x_{j_i}\|<\varepsilon/2\)。
因为
\[
 |(g-f)(y_i)|\leq |(g-f)(x_{j_i})|+2R\|y_i-x_{j_i}\|,
\]
只需把有限个 \(|(g-f)(x_{j_i})|\) 控制在 \(\varepsilon/2\) 内。
这些条件又包含 \(f\) 的某个 \(d\)-球：在距离和中分别保留第 \(j_i\) 项即可。
两种邻域系统因此等价。

由 Banach--Alaoglu 定理，这个度量空间紧，故序列紧。
也可以直接看到抽取过程：先逐个在有界标量列 \((f_n(x_j))_n\) 中抽子列，
再取对角子列。上述逼近估计保证这条子列在每个 \(x\in X\) 上取值都为 Cauchy 列。
令 \(f(x)\) 为逐点极限，线性关系可以逐点取极限，且 \(|f(x)|\leq R\|x\|\)。
于是 \(f\in RB_{X^*}\)，抽出的子列弱-*收敛到 \(f\)。
\end{proof}

范数有界在这里不可去掉：稠密集上的信息要推广到整个空间，靠的是对所有泛函共同有效的 \(2R\) 估计。
这个证明也说明，抽子列时测试函数不必事先包含空间中的每一个向量。
先在一个可数稠密族上安排收敛，再用一致估计补足其余测试，是后面弱解紧性论证的常用步骤。

不可分情形确实可能没有这样的子列。令 \(I=\mathcal P(\mathbb N)\)，
取 \(X=\ell^1(I)\)，即所有至多可数个非零坐标、绝对可和的族 \(x=(x_A)_{A\in I}\)，
范数为 \(\sum_A|x_A|\)。定义
\[
 f_n(x)=\sum_{A\in I}\mathbf1_A(n)x_A.
\]
则 \(\|f_n\|\leq1\)。任取子列 \((f_{n_k})\)，令 \(A=\{n_{2j}:j\geq1\}\)，
并以第 \(A\) 个坐标的单位向量 \(e_A\) 测试。
序列 \(f_{n_k}(e_A)\) 在零与一之间交替，不收敛。
所以这一对偶单位球虽然弱-*紧，却含有无弱-*收敛子列的序列。
一般紧性与序列紧性之间的区别不能仅凭记号中的收敛箭头消除。

% !TeX root = ../../Book1.tex
\section{反身性与弱紧性}
\label{b1:sec:0904}

% TODO: 撰写本节正文。

\subsection{反身空间}
\label{b1:subsec:090401}

待填充。

\subsection{有界序列的弱收敛子列}
\label{b1:subsec:090402}

待填充。

\subsection{Hilbert 空间中的弱紧性}
\label{b1:subsec:090403}

待填充。

\OptionalSubsection{Eberlein–Šmulian 定理}{b1:subsec:090404}

待填充。

% !TeX root = ../../Book1.tex
\section{凸性与弱下半连续性}
\label{b1:sec:0905}

上一节保证了自反空间中的有界序列可以抽取弱收敛子列。但在极值问题中，找到弱极限还不够：极限必须满足原来的约束，目标泛函的值也必须受到控制。凸性在这两个问题上都有作用。它使范数闭集保持弱闭，又使一大类范数下半连续泛函保持弱下半连续。以下先从凸集的闭包入手，再把集合的结论用于泛函。

\subsection{凸集}
\label{b1:subsec:090501}

回顾第八章的约定，集合 \(C\) 凸是指它包含任意两点间的线段；在复空间中，线段的参数仍取实数。反复运用这个条件可知，若 \(x_1,\ldots,x_m\in C\)，则
\[
  \sum_{j=1}^{m}\lambda_jx_j\in C,
  \quad \lambda_j\geq0,\quad \sum_{j=1}^{m}\lambda_j=1.
\]
这样的和称为这些点的\term[tu-zuhe]{凸组合}[convex combination]。与线性组合不同，凸组合不允许负系数，且系数之和固定为一。

上述有限点结论可以归纳证明。若 \(\lambda_m=1\)，和就是 \(x_m\)；若 \(\lambda_m<1\)，则先对前 \(m-1\) 个点使用归一化系数 \(\lambda_j/(1-\lambda_m)\)，得到 \(z\in C\)，原来的和便等于 \((1-\lambda_m)z+\lambda_mx_m\)。因此，两点间的线段条件已经包含所有有限平均的稳定性。

\begin{definition}[凸包与闭凸包]
\label{b1:def:convex-hull}
集合 \(A\subset X\) 的\term[tubao]{凸包}[convex hull]记为 \(\operatorname{co}A\)，是 \(A\) 中有限多个点的所有凸组合组成的集合，约定 \(\operatorname{co}\varnothing=\varnothing\)。它的范数闭包称为 \(A\) 的\term[bi-tubao]{闭凸包}[closed convex hull]。
\end{definition}

凸包确实是包含 \(A\) 的最小凸集。它包含 \(A\)；两个有限凸组合再取凸组合，仍是原来有限多个点的凸组合，所以它是凸的；任何包含 \(A\) 的凸集又必须包含所有这些组合。范数闭包也保持凸性：若 \(x_n,y_n\in C\) 分别依范数趋于 \(x,y\)，则对固定 \(t\in[0,1]\)，
\[
  (1-t)x_n+ty_n\longrightarrow(1-t)x+ty.
\]
因此，闭凸包是包含 \(A\) 的最小范数闭凸集。这些基本事实也可参见 Fremlin 对凸集与凸包的整理\cite[第2卷，2A5E]{Fremlin2016Measure}。

例如，两个不同点 \(u,v\) 的凸包只是线段 \([u,v]\)，而不包括整条仿射直线。接下来使用凸包时，目的是对已知向量作受控的平均，并不需要它们张成的整个线性空间。

\subsection{闭凸集的弱闭性}
\label{b1:subsec:090502}

弱拓扑比范数拓扑粗，故一般集合的弱闭包可能更大。凸集却有特殊之处：连续线性泛函能够把集外的点与整个闭凸集分开。这个几何事实正好能构造出排除该凸集的弱邻域。

\begin{theorem}[凸集的两种闭包]
\label{b1:thm:convex-weak-closure}
设 \(X\) 是赋范空间，\(C\subset X\) 凸。则 \(C\) 范数闭当且仅当 \(C\) 弱闭。进一步，以 \(\overline A^{\,\|\cdot\|}\) 和 \(\overline A^{\,w}\) 分别表示范数闭包与弱闭包，有
\[
  \overline C^{\,\|\cdot\|}=\overline C^{\,w}.
\]
\end{theorem}

\begin{proof}
弱闭集必范数闭，因为弱开集都是范数开集。反过来，设 \(C\) 非空、范数闭且凸。对任意 \(x_0\notin C\)，由\cref{b1:thm:geometric-hahn-banach}，存在 \(f\in X^*\) 及 \(\alpha\in\mathbb R\)，使
\[
  \sup_{z\in C}\operatorname{Re}f(z)
  \leq\alpha<\operatorname{Re}f(x_0).
\]
集合 \(\{x:\operatorname{Re}f(x)>\alpha\}\) 是 \(x_0\) 的弱开邻域，且与 \(C\) 不交。因此 \(X\setminus C\) 弱开，\(C\) 弱闭。空集的情形显然成立。

对任意凸集 \(C\)，它的范数闭包仍然凸，由刚证的结论也是弱闭集。因此 \(\overline C^{\,w}\subset\overline C^{\,\|\cdot\|}\)。反向包含对所有集合都成立，因为弱拓扑较粗。
\end{proof}

这里无需假设 \(X\) 完备。真正起作用的是凸分离；Fremlin 在更一般的局部凸空间中也以分离定理说明这一点\cite[第4卷，4A4Eb、4A4Ed]{Fremlin2016Measure}。由此还能把弱收敛转化为一种范数逼近，不过必须允许作凸组合。

\term*[mazur-yinli]{Mazur 引理}[Mazur's lemma]给出精确的结论。

\begin{theorem}[Mazur 引理]
\label{b1:thm:mazur}
若赋范空间 \(X\) 中的序列 \(x_n\rightharpoonup x\)，则对每个 \(n\) 可以选取 \(N_n\geq n\) 及非负数 \(\lambda_{n,k}\)，使
\[
  \sum_{k=n}^{N_n}\lambda_{n,k}=1,
  \quad y_n=\sum_{k=n}^{N_n}\lambda_{n,k}x_k,
  \quad \|y_n-x\|<\frac1n.
\]
换言之，\(x\) 可以由每个尾部的凸组合依范数逼近。
\end{theorem}

\begin{proof}
固定 \(n\)。由弱收敛，\(x\) 的每个弱邻域都包含某个 \(x_k\)，其中 \(k\geq n\)，因而 \(x\) 属于 \(\{x_k:k\geq n\}\) 的弱闭包，更属于其凸包的弱闭包。对这个凸包应用\cref{b1:thm:convex-weak-closure}，便知 \(x\) 属于它的范数闭包。于是存在距离 \(x\) 小于 \(1/n\) 的有限凸组合。补入系数为零的项即可写成上述形式。
\end{proof}

以 \(\ell^2\) 中的标准单位向量 \(e_n\) 为例。由 Hilbert 空间的 Riesz 表示，每个连续线性泛函在 \(e_n\) 上的值趋于零，故 \(e_n\rightharpoonup0\)，但 \(\|e_n\|=1\)。它没有范数收敛于零的子列，尾部平均却满足
\[
  \left\|\frac1n\sum_{k=n}^{2n-1}e_k\right\|=\frac1{\sqrt n}
  \longrightarrow0.
\]
因此，Mazur 引理提供的是凸组合，而不是强收敛子列。

\subsection{范数的弱下半连续性}
\label{b1:subsec:090503}

极小值问题不要求泛函值沿极限保持相等。若极限点的值比逼近序列更小，反而有利；需要排除的是取极限时向上跳跃。

\begin{definition}[下半连续性]
\label{b1:def:lower-semicontinuity}
设 \(Y\) 是拓扑空间，\(J:Y\to(-\infty,+\infty]\)。若对每个 \(a\in\mathbb R\)，\term[zishui-pingji]{子水平集}[sublevel set]
\[
  \{y\in Y:J(y)\leq a\}
\]
都是闭集，则称 \(J\) 具有\term[xiaban-lianxuxing]{下半连续性}[lower semicontinuity]。在赋范空间上，相对于范数拓扑与弱拓扑的下半连续性分别称为强下半连续性与\term[ruo-xiaban-lianxuxing]{弱下半连续性}[weak lower semicontinuity]。
\end{definition}

这个定义等价于 \(\{y:J(y)>a\}\) 对所有实数 \(a\) 都开。特别地，若 \(J\) 下半连续且 \(y_n\to y\)，则
\begin{equation}
\label{b1:eq:lsc-sequence}
  J(y)\leq\liminf_{n\to\infty}J(y_n).
\end{equation}
事实上，对任意 \(a<J(y)\)，上述开集是 \(y\) 的邻域，故最终 \(J(y_n)>a\)；让 \(a\) 趋近 \(J(y)\) 即得结论，包括 \(J(y)=+\infty\) 的情形。

在度量空间中，序列条件也能反过来推出下半连续性。若某个子水平集不闭，取其闭包中不属于它的点 \(y\)，并在每个半径 \(1/n\) 的球内取该子水平集中的点 \(y_n\)。此时 \(y_n\to y\)，但 \(J(y_n)\leq a<J(y)\)，与\eqref{b1:eq:lsc-sequence}矛盾。这就给出了强下半连续性的序列判别。对于一般弱拓扑，不能未经证明便把这个序列判别当成定义：上述逆推依赖于度量空间的球邻域基。以下仍用子水平集的弱闭性证明弱下半连续性，而只在取弱收敛子列时使用其必然成立的序列不等式。一般拓扑下的相应性质可参见\cite[第4卷，4A2Bd]{Fremlin2016Measure}。

\begin{proposition}[范数的弱下半连续性]
\label{b1:prop:norm-weak-lsc}
赋范空间 \(X\) 的范数是弱下半连续函数。因而若 \(x_n\rightharpoonup x\)，则
\[
  \|x\|\leq\liminf_{n\to\infty}\|x_n\|.
\]
\end{proposition}

\begin{proof}
由\cref{b1:cor:norming-functional}，
\[
  \|x\|=\sup_{\substack{f\in X^*\\\|f\|\leq1}}|f(x)|.
\]
对 \(r\geq0\)，相应子水平集可写为
\[
  \{x:\|x\|\leq r\}
  =\bigcap_{\substack{f\in X^*\\\|f\|\leq1}}
    \{x:|f(x)|\leq r\}.
\]
每个 \(f\) 弱连续，右侧各集合弱闭，其交也弱闭；当 \(r<0\) 时子水平集为空。序列不等式于是来自\eqref{b1:eq:lsc-sequence}。
\end{proof}

注意，范数的弱下半连续性并不等于弱连续性。上面的 \(e_n\rightharpoonup0\) 正说明范数可以从恒为一降到零。

\subsection{凸泛函的弱下半连续性}
\label{b1:subsec:090504}

为了把约束也纳入泛函，我们允许它取 \(+\infty\)，但不允许取 \(-\infty\)。约定 \(J:X\to(-\infty,+\infty]\) 不恒等于 \(+\infty\)，即其\term[youxiao-dingyiyu]{有效定义域}[effective domain]
\[
  \operatorname{dom}J=\{x\in X:J(x)<+\infty\}
\]
非空。

\begin{definition}[凸泛函与严格凸泛函]
\label{b1:def:convex-functional}
若对所有 \(x,y\in X\) 及 \(0<t<1\)，有
\[
  J\bigl((1-t)x+ty\bigr)\leq(1-t)J(x)+tJ(y),
\]
则称 \(J\) 为\term[tu-fanhan]{凸泛函}[convex functional]。若当 \(x,y\in\operatorname{dom}J\) 且 \(x\ne y\) 时总为严格不等式，则称其为\term[yange-tu-fanhan]{严格凸泛函}[strictly convex functional]。
\end{definition}

取 \(0<t<1\) 避免了端点处 \(0\cdot(+\infty)\) 的约定。凸泛函的有效定义域是凸集，每个子水平集也是凸集：若两点的值均不超过 \(a\)，其凸组合的值便不超过 \((1-t)a+ta=a\)。因此，集合的闭性定理可以直接用于泛函。

这些性质还可以用一个集合统一表示：在 \(X\times\mathbb R\) 中令
\[
  G_J=\{(x,r):J(x)\leq r\}.
\]
\(J\) 凸当且仅当 \(G_J\) 凸。一个方向由凸不等式直接给出；另一个方向，对 \(J(x),J(y)\) 都有限的点，把 \(\bigl(x,J(x)\bigr)\) 与 \(\bigl(y,J(y)\bigr)\) 作凸组合即可。若其中一个值为无穷，凸不等式本来就成立。

相对于 \(X\) 上的指定拓扑与 \(\mathbb R\) 的通常拓扑，\(J\) 下半连续也等价于 \(G_J\) 闭。若 \(J\) 下半连续而 \(r<J(x)\)，选择 \(r<a<J(x)\)，则
\[
  \{y:J(y)>a\}\times(-\infty,a)
\]
是 \((x,r)\) 的开邻域，并且避开 \(G_J\)，故其补集开。反之，若 \(G_J\) 闭，则固定 \(a\) 后，连续映射 \(x\mapsto(x,a)\) 下的逆像 \(\{x:J(x)\leq a\}\) 闭。这样看来，凸性与下半连续性分别表达了同一个集合的凸性与闭性。不过下面的证明只需处理子水平集，不必另行研究乘积空间。

\begin{theorem}[凸泛函的两种下半连续性]
\label{b1:thm:convex-weak-lsc}
赋范空间上的凸泛函 \(J:X\to(-\infty,+\infty]\) 强下半连续当且仅当弱下半连续。
\end{theorem}

\begin{proof}
若 \(J\) 强下半连续，则每个子水平集范数闭，又因凸性而由\cref{b1:thm:convex-weak-closure}知其弱闭。这正是弱下半连续性的定义。反之，弱闭集必范数闭，所以弱下半连续总蕴含强下半连续，这个方向甚至不需要凸性。
\end{proof}

一个常用的做法是给非空约束集 \(C\) 配上泛函
\[
  \iota_C(x)=
  \begin{cases}
    0,&x\in C,\\
    +\infty,&x\notin C.
  \end{cases}
\]
它与测度论中的特征函数不同：集外的值是 \(+\infty\)。当 \(C\) 凸时 \(\iota_C\) 凸；当 \(C\) 范数闭时 \(\iota_C\) 强下半连续，因为其子水平集只有空集和 \(C\)。因此，闭凸约束可以用弱下半连续泛函表达。把能量 \(J\) 在 \(C\) 上最小化，也就等同于把 \(J+\iota_C\) 在全空间上最小化。

凸性不能仅凭范数连续性替代。函数 \(J(x)=-\|x\|^2\) 在 \(\ell^2\) 上范数连续，但对 \(e_n\rightharpoonup0\)，有
\[
  J(0)=0>-1=\liminf_{n\to\infty}J(e_n),
\]
所以它并不弱下半连续。

\subsection{变分法中的意义}
\label{b1:subsec:090505}

现在考虑
\[
  m=\inf_{x\in C}J(x).
\]
求解这一问题的第一步通常不是猜出满足某个方程的向量，而是选取使能量趋于 \(m\) 的序列，再设法保留它的极限。由此得到的存在性证明称为\term[bianfen-zhijiefa]{变分直接法}[direct method in the calculus of variations]。

\begin{theorem}[变分直接法]
\label{b1:thm:direct-method}
设 \(X\) 是自反 Banach 空间，\(C\subset X\) 非空且弱闭。设 \(J:C\to(-\infty,+\infty]\) 在 \(C\) 上有下界、至少在一点取有限值，并相对于 \(C\) 的弱子空间拓扑下半连续。再设 \(J\) 满足增长条件
\begin{equation}
\label{b1:eq:direct-method-growth}
  J(x)\longrightarrow+\infty
  \quad\text{当 }x\in C,\ \|x\|\longrightarrow\infty.
\end{equation}
则存在 \(u\in C\)，使 \(J(u)=\inf_{x\in C}J(x)\)。若另有 \(C\) 凸且 \(J\) 严格凸，则此极小点唯一。
\end{theorem}

\begin{proof}
由下界和有限值假设，\(m=\inf_CJ\) 是实数。对每个 \(n\)，选取 \(x_n\in C\)，使
\[
  m\leq J(x_n)<m+\frac1n.
\]
增长条件等价于每个有限子水平集都有界：对给定 \(a\in\mathbb R\)，存在 \(R\)，使 \(x\in C\) 且 \(\|x\|>R\) 时 \(J(x)>a\)。取 \(a=m+1\)，便知 \((x_n)\) 有界。

由\cref{b1:thm:reflexive-weak-subsequence}，存在子列 \(x_{n_k}\rightharpoonup u\)。由于 \(C\) 弱闭，\(u\in C\)；由于 \(J\) 弱下半连续，
\[
  J(u)\leq\liminf_{k\to\infty}J(x_{n_k})=m.
\]
而 \(m\) 是下确界，所以 \(J(u)=m\)。

若 \(u,v\) 是两个不同的极小点，且 \(C\) 凸、\(J\) 严格凸，则 \((u+v)/2\in C\)，且
\[
  J\left(\frac{u+v}{2}\right)
  <\frac{J(u)+J(v)}2=m,
\]
与下确界的定义矛盾。
\end{proof}

证明实际上只需要某个 \(a>m\) 对应的子水平集有界，因为极小化序列最终落在其中。下界假设则确保 \(m\ne-\infty\)，不能在选取 \(m+1/n\) 时悄悄省掉。自反性负责提供弱收敛子列，弱闭约束负责保留可行性，下半连续性负责比较能量；增长条件的作用只是把极小化序列限制在有界区域内。

以 Hilbert 空间 \(H\) 为例，给定 \(\ell\in H^*\) 及非空闭凸集 \(C\)，考虑能量
\[
  E(x)=\frac12\|x\|^2-\operatorname{Re}\ell(x),
  \quad x\in C.
\]
它是范数连续的，并满足
\[
  E\bigl((1-t)x+ty\bigr)
  =(1-t)E(x)+tE(y)-\frac{t(1-t)}2\|x-y\|^2,
\]
故严格凸。由\cref{b1:thm:convex-weak-lsc}，\(E\) 弱下半连续。又有
\[
  E(x)\geq\frac12\|x\|^2-\|\ell\|\cdot\|x\|
  =\frac12(\|x\|-\|\ell\|)^2-\frac12\|\ell\|^2,
\]
所以它有下界，并满足\eqref{b1:eq:direct-method-growth}。由\cref{b1:cor:hilbert-reflexive}，\(H\) 自反；由闭凸集的弱闭性和直接法，\(E\) 在 \(C\) 上有唯一极小点。

这一结论也能与第八章的投影问题直接对应。写 \(\ell(x)=\langle x,h\rangle\)，则
\[
  E(x)=\frac12\|x-h\|^2-\frac12\|h\|^2.
\]
故极小点就是 \(h\) 在 \(C\) 中的最近点，与\cref{b1:thm:hilbert-convex-minimum}一致。直接法不必预先解出这个点；它先由空间结构与能量性质保证点的存在，随后才可研究其特征方程或近似计算。

具体地，对极小点 \(u\) 和任意 \(v\in C\)，凸性保证 \(u+t(v-u)\in C\)。展开能量差得
\[
  0\leq E\bigl(u+t(v-u)\bigr)-E(u)
  =t\operatorname{Re}\langle u-h,v-u\rangle
    +\frac{t^2}{2}\|v-u\|^2.
\]
除以 \(t>0\) 再令 \(t\downarrow0\)，得到
\[
  \operatorname{Re}\langle u-h,v-u\rangle\geq0
  \quad(v\in C).
\]
反过来，这个不等式与能量差公式在 \(t=1\) 时的表达一起，又能推出 \(E(v)\geq E(u)\)。所以它是极小点的充要条件。无约束时可以取 \(v=h\)，从而得到 \(u=h\)；一般凸约束只允许沿留在 \(C\) 内的方向变化，得到的便是不等式而非任意方向上的等式。

即使在一维，若缺少限制极小化序列的条件，结论也可能失败。在 \(C=\mathbb R\) 上，\(J(t)=e^{-t}\) 连续、严格凸且下界为零，却没有极小点：能量趋于零的序列必向正无穷远离。若约束本身不闭，也有另一种失败方式：\(J(t)=t^2\) 在 \(C=(0,+\infty)\) 上满足上述增长条件，却只能在不属于 \(C\) 的边界点达到下确界。它们分别说明，控制序列的位置与保留极限点的约束是存在性证明中不同的要求。

% 本章小结（不计入编号）
\begin{chaptersummary}
弱收敛测试 \(X^*\) 的全部成员，弱-*\WeakStarJoin{}收敛则在一个已经写成对偶的空间上只测试原空间。
两者都把无穷维收敛拆成标量信息，但标量信息不能自动恢复范数，也不能无条件用序列描述整个拓扑。
弱收敛序列必有界。当原空间为 Banach 空间时，一致有界原理保证弱-*\WeakStarJoin{}收敛序列范数有界；若原空间不完备，这一一般结论可能失效。

\ProseParagraphBreak

Banach--Alaoglu 定理将对偶单位球放入紧乘积，再以线性关系截出闭子集。
可分性把有界球上的全部测试压缩到可数稠密族，因而得到度量和收敛子列。
Goldstine 定理与自然嵌入进一步把自反性刻画为原空间闭单位球的弱紧性；
自反空间的任意有界序列都有弱收敛子列，不要求整个空间可分。

\ProseParagraphBreak

凸性补偿了拓扑减弱造成的一部分信息损失：凸集的范数闭包与弱闭包相同，
弱收敛序列的尾部凸组合能够范数逼近极限，凸的范数下半连续泛函也是弱下半连续的。
这些性质使极小化列的弱极限仍满足约束，且泛函值不超过原列的下极限。
下一章将在 \(L^p\) 空间中落实对偶与自反性，说明上述方法对哪些函数空间有效，
以及 \(p=1\) 和 \(p=\infty\) 为什么必须分别处理。
\end{chaptersummary}


\BookExercises

% 以下十一题由本章工具组织并给出独立解答，不逐字摘录文献习题。

\begin{exercise}\label{b1:ex:09-basic-convergence}
分别判断下列序列的强收敛、弱收敛或弱-*\WeakStarJoin{}收敛，并说明所断言的较强收敛为何成立或失败。
\begin{enumerate}
\item \(\ell^2\) 中的 \(x^{(n)}=e_n/n\)；
\item \(\ell^2\) 中的 \(y^{(n)}=e_n\)；
\item 把 \(\ell^\infty\) 识别为 \((\ell^1)^*\)，并令
\[
 z^{(n)}=(\underbrace{1,\ldots,1}_{n\text{ 项}},0,0,\ldots).
\]
\end{enumerate}
\end{exercise}
\begin{solution}
第一列满足 \(\|x^{(n)}\|_2=1/n\to0\)，故强收敛于零，也弱收敛于零。

第二列的范数恒为一，因而不强收敛于零。对每个 \(h\in\ell^2\)，有
\(\langle e_n,h\rangle=\overline{h_n}\to0\)；由 Riesz 表示定理，\(y^{(n)}\rightharpoonup0\)。

第三列范数恒为一，每个固定坐标最终等于一。由\cref{b1:cor:linfty-weak-star-coordinate}，
\[
 z^{(n)}\overset{*}{\rightharpoonup}\boldsymbol 1=(1,1,\ldots).
\]
但 \(\|z^{(n)}-\boldsymbol1\|_\infty=1\) 对每个 \(n\) 都成立，所以它不强收敛于这个弱-*\WeakStarJoin{}极限。
\end{solution}

\begin{exercise}\label{b1:ex:09-basic-weak-neighborhood}
设 \(X\) 是赋范空间，\(x_0\in X\)，\(f_1,\ldots,f_m\in X^*\)，且 \(\varepsilon_1,\ldots,\varepsilon_m>0\)。令
\[
 W=\{x\in X:|f_j(x-x_0)|<\varepsilon_j\text{ 对 }1\leq j\leq m\}.
\]
证明 \(W\) 是 \(x_0\) 的弱邻域，并写出一个包含于 \(W\) 的标准基本邻域 \(U(x_0;F,\varepsilon)\)。若 \(X\) 无穷维，再证明 \(W\) 包含一条经过 \(x_0\) 的仿射直线，因而在范数下无界。
\end{exercise}
\begin{solution}
取 \(F=\{f_1,\ldots,f_m\}\) 及 \(\varepsilon=\min_j\varepsilon_j\)，则
\[
 U(x_0;F,\varepsilon)\subset W,
\]
所以 \(W\) 是弱邻域。若 \(X\) 无穷维，线性映射
\[
 T:X\to\mathbb K^m,\qquad Tx=(f_1(x),\ldots,f_m(x))
\]
不可能单射。取非零 \(v\in\ker T\)，便有 \(f_j(tv)=0\) 对所有标量 \(t\) 和全部 \(j\) 成立。因此 \(x_0+tv\in W\) 对每个 \(t\) 成立，而这条直线在范数下无界。
\end{solution}

\begin{exercise}\label{b1:ex:09-weak-star-coordinate-test}
设 \(u^{(n)}=(u_k^{(n)})\in\ell^\infty\)，并假设
\[
 \sup_n\|u^{(n)}\|_\infty\leq M,
 \qquad u_k^{(n)}\longrightarrow u_k\quad(k\geq1).
\]
不用直接引用坐标判别推论，证明 \(u=(u_k)\in\ell^\infty\)，并从 \(\ell^1\) 配对的首项与尾项估计推出 \(u^{(n)}\overset{*}{\rightharpoonup}u\)。再用 \(v^{(n)}=ne_n\) 说明：若删去统一范数界，逐坐标收敛不足以推出弱-*\WeakStarJoin{}收敛。
\end{exercise}
\begin{solution}
逐坐标取极限给出 \(|u_k|\leq M\)，故 \(u\in\ell^\infty\)。对任意 \(a\in\ell^1\) 及正整数 \(N\)，有
\[
 \left|\sum_{k=1}^{\infty}a_k(u_k^{(n)}-u_k)\right|
 \leq\sum_{k=1}^N|a_k|\,|u_k^{(n)}-u_k|
     +2M\sum_{k>N}|a_k|.
\]
先选 \(N\) 控制尾和，再对有限个首项取 \(n\to\infty\)，便得到全部 \(\ell^1\) 测试下的收敛。

序列 \(v^{(n)}=ne_n\) 的每个固定坐标都最终为零。令 \(a_{2^j}=2^{-j}\)，其余坐标为零，则 \(a\in\ell^1\)，而
\[
 \sum_k a_kv_k^{(n)}=na_n
\]
在 \(n=2^j\) 时等于一，在其余无限多个 \(n\) 上等于零，故不收敛于零。于是 \(v^{(n)}\) 并不弱-*\WeakStarJoin{}收敛于零。
\end{solution}

\begin{exercise}\label{b1:ex:09-varying-tests}
设 \(X\) 是 Banach 空间。若 \(x_n\rightharpoonup x\) 且 \(f_n\to f\) 于 \(X^*\) 的范数，
证明 \(f_n(x_n)\to f(x)\)。若改为 \(x_n\to x\) 于范数而 \(f_n\) 弱-*\WeakStarJoin{}趋于 \(f\)，结论是否仍成立？
最后说明同时只要求 \(x_n\rightharpoonup x\) 与 \(f_n\) 弱-*\WeakStarJoin{}趋于 \(f\) 时，结论可以失败。
\end{exercise}
\begin{solution}
第一种情形中 \(M=\sup_n\|x_n\|<\infty\)，且
\[
 |f_n(x_n)-f(x)|\leq M\|f_n-f\|+|f(x_n-x)|\longrightarrow0.
\]
第二种情形中，一致有界原理给出 \(M'=\sup_n\|f_n\|<\infty\)，于是
\[
 |f_n(x_n)-f(x)|\leq M'\|x_n-x\|+|f_n(x)-f(x)|\longrightarrow0.
\]
若两边都只有较弱的收敛，取 \(X=\ell^2\)、\(x_n=e_n\)、\(f_n(y)=y_n\)。
则 \(x_n\rightharpoonup0\)，每个固定 \(y\in\ell^2\) 的坐标 \(y_n\to0\)，故 \(f_n\) 弱-*\WeakStarJoin{}趋于零；
但 \(f_n(x_n)=1\)。逐个固定测试的极限，不能直接用于随 \(n\) 一同改变的测试对象。
\end{solution}

\begin{exercise}\label{b1:ex:09-sphere-closure}
设 \(H\) 是无穷维 Hilbert 空间。证明其单位球面 \(S_H=\{x:\|x\|=1\}\) 的弱闭包是闭单位球。
并证明闭球中的每一点都可以由球面中的一条序列弱逼近。
\end{exercise}
\begin{solution}
范数弱下半连续，故 \(B_H\) 弱闭，得到 \(\overline{S_H}^{\,w}\subset B_H\)。
固定 \(x\in B_H\)。若 \(\|x\|=1\)，取常值列。
若 \(\|x\|<1\)，则 \(x^\perp\) 仍无穷维，可以逐次选出标准正交序列 \((e_n)\subset x^\perp\)。
由 Bessel 不等式和 Riesz 表示，\(e_n\rightharpoonup0\)。于是
\[
 y_n=x+\sqrt{1-\|x\|^2}\,e_n
\]
满足 \(\|y_n\|=1\) 且 \(y_n\rightharpoonup x\)。这既给出反向包含，也给出所求序列。
球面范数闭而不弱闭；它不是凸集，故不与凸集的闭包定理冲突。
\end{solution}

\begin{exercise}\label{b1:ex:09-weak-cauchy-reflexive}
证明自反 Banach 空间中的每个弱 Cauchy 序列都有弱极限。
指出证明的哪一步不适用于 \(c_0\) 中前 \(n\) 项为一、其余为零的序列。
\end{exercise}
\begin{solution}
设 \((x_n)\) 为弱 Cauchy 列。对每个 \(f\in X^*\)，标量列 \(\bigl(f(x_n)\bigr)\) 有界。
对 Banach 空间 \(X^*\) 上的泛函族 \(\{J_Xx_n\}\) 使用一致有界原理，得 \(\sup_n\|x_n\|<\infty\)。
自反性给出子列 \(x_{n_k}\rightharpoonup x\in X\)。对每个 \(f\)，
原标量列是 Cauchy 列，又有子列趋于 \(f(x)\)，故全列也趋于 \(f(x)\)。因此 \(x_n\rightharpoonup x\)。
在所给 \(c_0\) 例子中，任何弱极限都必须逐坐标等于一，这个向量不在 \(c_0\) 中。
该空间不自反，有界性不能提供所需的弱收敛子列。
\end{solution}

\begin{exercise}\label{b1:ex:09-continuous-evaluation}
设 \(X\) 是赋范空间，\(L:X^*\to\mathbb K\) 为线性泛函。
证明 \(L\) 对 \(\sigma(X^*,X)\) 连续，当且仅当存在 \(x\in X\)，使 \(L(f)=f(x)\) 对所有 \(f\in X^*\) 成立。
\end{exercise}
\begin{solution}
逐点评值按弱-*\WeakStarJoin{}拓扑的定义连续，只需证必要性。
由零点连续性，可找 \(x_1,\ldots,x_m\in X\) 及 \(\delta>0\)，使
\[
 \max_i|f(x_i)|<\delta\quad\Longrightarrow\quad |L(f)|<1.
\]
若 \(f(x_i)=0\) 对所有 \(i\) 成立，任意标量倍 \(tf\) 都满足左边条件，
故 \(|tL(f)|<1\) 对任意 \(t\) 成立，推出 \(L(f)=0\)。
令 \(Tf=\bigl(f(x_1),\ldots,f(x_m)\bigr)\)。上述结论表明 \(L\) 在 \(\ker T\) 上为零，
所以 \(A(Tf)=L(f)\) 定义了 \(T(X^*)\) 上的线性泛函。
将 \(A\) 代数延拓到有限维空间 \(\mathbb K^m\)，写成 \(A(z)=\sum_i a_iz_i\)。
于是 \(L(f)=\sum_i a_if(x_i)=f(\sum_i a_ix_i)\)，取 \(x=\sum_i a_ix_i\) 即可。
这个证明说明，一个弱-*\WeakStarJoin{}连续的线性测试最终只需有限个原空间向量，合并后就是单点评值。
\end{solution}

\begin{optionalexercise}\label{b1:ex:09-compact-improves}
设 \(X,Y\) 是 Banach 空间，\(T:X\to Y\) 是紧线性算子。
证明 \(x_n\rightharpoonup x\) 蕴含 \(Tx_n\to Tx\) 于范数。
\end{optionalexercise}
\begin{solution}
任意 \(g\in Y^*\) 与 \(T\) 的复合属于 \(X^*\)，故 \(Tx_n\rightharpoonup Tx\)。
同时 \((x_n)\) 范数有界，紧性保证 \((Tx_n)\) 的任意子列都有范数收敛子子列。
假如 \(Tx_n\) 不范数趋于 \(Tx\)，就有 \(\varepsilon>0\) 及子列满足
\(\|Tx_{n_k}-Tx\|\geq\varepsilon\)。再抽范数收敛子子列，设其极限为 \(y\)。
范数收敛蕴含弱收敛，而原弱极限为 \(Tx\)，由弱极限唯一性得 \(y=Tx\)，矛盾。
紧性在这里把原本只有测试意义的收敛提升为统一的范数控制。
\end{solution}

\begin{exercise}\label{b1:ex:09-double-annihilator}
设 \(A\) 是赋范空间 \(X\) 的任意子集，记
\[
 A^\circ=\{f\in X^*:f(a)=0\text{ 对每个 }a\in A\}.
\]
证明 \(A^\circ=(\operatorname{span}A)^\circ\)，并利用\cref{b1:prop:double-annihilator}推出
\[
 \overline{\operatorname{span}A}^{\,\|\cdot\|}
 =\overline{\operatorname{span}A}^{\,w}
 =\{x\in X:f(x)=0\text{ 对所有 }f\in A^\circ\}.
\]
特别地，证明 \(\operatorname{span}A\) 在 \(X\) 中稠密，当且仅当 \(A^\circ=\{0\}\)。
\end{exercise}
\begin{solution}
泛函若在 \(A\) 上为零，由线性性就在 \(\operatorname{span}A\) 上为零；反向包含显然。因此
\(A^\circ=(\operatorname{span}A)^\circ\)。把 \(M=\operatorname{span}A\) 代入\cref{b1:prop:double-annihilator}，再用这个零化子等式，就得到题中的三个集合相等。

若 \(\operatorname{span}A\) 稠密，连续泛函在 \(A\) 上为零便在整个 \(X\) 上为零，所以 \(A^\circ=\{0\}\)。反过来，若 \(A^\circ=\{0\}\)，上述等式最右侧没有任何非零测试条件，等于整个 \(X\)，故 \(\overline{\operatorname{span}A}=X\)。
\end{solution}

\begin{exercise}\label{b1:ex:09-norm-attainment}
设 \(X\ne\{0\}\) 是自反 Banach 空间。证明每个 \(f\in X^*\) 都在闭单位球上取得范数。
再证明 \(c_0\) 上的泛函 \(f(x)=\sum_{n\geq1}2^{-n}x_n\) 范数等于一，
但没有 \(x\in B_{c_0}\) 使 \(|f(x)|=1\)。
\end{exercise}
\begin{solution}
闭单位球弱紧，\(x\mapsto|f(x)|\) 弱连续，故它在球上取得最大值，按范数定义此值就是 \(\|f\|\)。
在第二个例子中，绝对收敛及 \(|f(x)|\leq\|x\|_\infty\) 给出 \(\|f\|\leq1\)。
以头 \(N\) 项为一的有限支撑向量测试，得到 \(\|f\|\geq\sum_{n=1}^N2^{-n}\)，令 \(N\to\infty\) 得到等号。
但对任意 \(x\in B_{c_0}\)，由 \(x_n\to0\) 可选一个 \(j\) 使 \(|x_j|<1\)，从而
\[
 |f(x)|\leq\sum_n2^{-n}|x_n|
 \leq1-2^{-j}(1-|x_j|)<1.
\]
因此对偶范数中的上确界，在一般 Banach 空间中不一定是最大值。
\end{solution}

\begin{exercise}\label{b1:ex:09-noncoercive-quadratic}
在实 Hilbert 空间 \(\ell^2\) 中，令
\[
 C=\left\{x:\sum_{n=1}^{\infty}\frac{x_n}{n}=1\right\},
 \quad F(x)=\sum_{n=1}^{\infty}\frac{x_n^2}{n^2}.
\]
证明 \(C\) 是非空弱闭凸集，\(F\) 是凸的弱下半连续泛函，
但 \(F\) 在 \(C\) 上没有极小点。找出直接法中失效的步骤。
\end{exercise}
\begin{solution}
向量 \((1/n)\in\ell^2\)，故约束是一个连续线性泛函的水平集，弱闭且凸；\(e_1\in C\)，所以非空。
映射 \(T(x_n)=(x_n/n)\) 有界，\(F(x)=\|Tx\|^2\) 范数连续且凸，因而弱下半连续。
也可把 \(F\) 写成有限坐标连续凸函数之上确界，直接得到同一结论。

令 \(x_n^{(N)}=n/N\) 当 \(1\leq n\leq N\)，其余坐标为零。
则 \(x^{(N)}\in C\)，而 \(F(x^{(N)})=1/N\)，故 \(\inf_CF=0\)。
若 \(F(x)=0\)，每个坐标都为零，与约束矛盾，所以该下确界不取到。
这里 \(\|x^{(N)}\|^2=N^{-2}\sum_{n=1}^N n^2\to\infty\)，
极小化列没有从能量得到范数有界性。更一般地，任何极小化列若有有界子列，
自反性、弱闭性和弱下半连续性便会给出极小点，与刚才的结论矛盾。
所以有下界不等于有足够控制无穷远行为的增长条件。
\end{solution}
